\section{Discussion}
\label{sec:discussion}

For Markov sequences, the structure a denoiser can use is not in the score, which diffusion makes non-local, but in the
posterior, which stays a tree. BP turns a kernel, a function of two numbers, into the exact score at every noise level,
and EM learns the kernel from noisy data with BP as its E-step, without the need to differentiate through it. A small neural
kernel removes the need to choose a parametric family to estimate the Kernel, and remains accurate and more efficient than score denoisers.

\paragraph{What this suggests for diffusion models.} If the data is first-order Markov, do not learn the score; learn the kernel and use BP to get the score.

\paragraph{Limits.} The Markov order has to be known: on pairs of sites a grid step costs $O(N_g^3)$ instead of
$O(N_g^2)$. Each site carries one number here; a state with $d$ numbers needs a grid of $N_g^d$ points. The graph must be
a tree, and the grid limits how small $t$ can be. Finally, a large network closes the gap when data are plentiful, so the
gain we measure is in data efficiency, which matters most where scientific data are scarce.

\paragraph{Related work.} EM with forward--backward as the E-step is the classical way to fit hidden Markov and
state-space models \cite{baum1970,rabiner1989tutorial,shumway1982em,cappe2005}, and grids for non-Gaussian filtering go
back to \cite{kitagawa1987}; here the hidden chain is the clean data, the emission is the diffusion noise, and the
fitted model is used for its score. Diffusion models can also be trained from corrupted data
\cite{daras2023ambient,rozet2024learning,bai2024em}, but they learn a network score. Analytic and structured denoisers
\cite{kamb2024analytic,kadkhodaie2024generalization}, belief propagation on hierarchical data
\cite{garnierbrun2024transformers,sclocchi2025phase,mei2024unets} and the statistical physics of diffusion
\cite{biroli2023generative,biroli2024dynamical,achilli2025memorization,bonnaire2025memorize,garnierbrun2026biased,garnierbrun2026foam}
ask what structure a denoiser exploits. Markov sequences are the simplest (possibly) non-Gaussian case in which the exact (up to discretization and kernel estimation error) denoiser
is computable and can be learned from noisy data alone.
